\section{Natural linearization and cohomological input}\label{sec:cohomological-input}
Write $n=\dim X$ and $\chi^T(X,\mathcal O_X)=\sum_q(-1)^q\dim H^q(X,\mathcal O_X)^T$. The transfer theorem permits a chosen character. To obtain its input from cohomology, we now use the natural anticanonical character. The following result is the exact input from the index companion \cite[Invariant anticanonical index theorem]{CompanionI}.

\begin{proposition}[Invariant anticanonical index]\label{har:prop:index}
Let $F$ be a compact complex manifold with a holomorphic action of a compact real torus $T$. Give $D=-K_F$ its natural linearization. There are an integer $a>0$ and a polynomial $P\in\Q[t]$ such that
\[
 \sum_q(-1)^q\dim H^q(F,maD)^T=P(m)\quad(m>0),\qquad
 P(0)=\sum_q(-1)^q\dim H^q(F,\mathcal O_F)^T.
\]
The statement includes the trivial torus, an empty fixed locus, and zero-dimensional fixed components; it assumes no positivity, projectivity, or K\"ahler condition.
\end{proposition}

In particular the value at zero is the actual index, not a formal extrapolation of an eventual polynomial. We use this theorem to select a fixed cohomological degree on an unbounded sequence.

\begin{lemma}[Forms and determinant sections]\label{har:lem:forms}
Let \(X\) be smooth connected projective, let \(L=-K_X\) have a smooth
semipositive metric, and let an algebraic torus \(T\) act on \(X\).
If \(\chi^T(X,\mathcal O_X)\ne0\), then for some fixed \(q\) there are
nonzero invariant sections of
\(\Omega_X^{n-q}\otimes L^{m+1}\) for unbounded positive integers \(m\).
If these do not already give a nonzero invariant section of a positive power of $L$, there
is an integer $e>0$, a linearized line bundle \(M\), a nonzero map
\(M\to\bigwedge^e\Omega_X^{n-q}\), and invariant sections
\[
 0\ne s_i\in H^0(X,M+a_iL)^T,\qquad a_i\longrightarrow\infty.
\]
\end{lemma}

\begin{proof}
The polynomial in Proposition~\ref{har:prop:index} is not zero. A fixed
cohomological degree therefore has nonzero invariant cohomology along an
unbounded sequence. Average a Kähler form and the metric weights over
the maximal compact torus, and denote the resulting curvature of \(L\)
by \(\theta\). The smooth semipositive hard Lefschetz theorem
\cite[Corollary~2.1.2]{DPS2001}, applied to \(mL=K_X+(m+1)L\),
gives the asserted forms. Concretely, if \(u\) is harmonic
of type \((n,q)\) with values in \((m+1)L\), Bochner--Kodaira gives
\[
 0=\|(D')^*u\|^2+\langle[(m+1)\theta,\Lambda]u,u\rangle.
\]
Both terms are nonnegative. Applying the complex-linear star on form
indices to $u$ therefore gives a holomorphic form of type \((n-q,0)\).
The construction is equivariant.
For \(q=0\) it already yields invariant sections of \(mL\); for \(q=n\)
the holomorphic form itself is a section of \((m+1)L\). If
\(r=\max\operatorname{rank}\theta=n\) and
\(q>0\), the curvature term makes this form vanish on a nonempty open set,
hence everywhere, so this case cannot occur.

The determinant extraction follows the method of
\cite[Lemma~5.1]{LMPTX2023} and \cite[Lemma~4.1]{LP2018}.
We spell out the unbounded-twist step, which here also applies when
$L$ is numerically trivial. Choose a rational frame of $L$ and divide
each form by the corresponding power of that frame. Their span over
$\C(X)$ is a subspace of the generic fiber of $\Omega_X^{n-q}$,
of some rank $e$. Changing the frame rescales each vector by a nonzero
rational function, so this subspace is intrinsic. Saturate its top
exterior power in $\bigwedge^e\Omega_X^{n-q}$ and call the resulting
rank-one subsheaf $M$. It is reflexive on smooth $X$, hence a line
bundle. The generic span is torus-stable, so its saturation carries the
induced linearization.

Choose finitely many of the original forms as a generic basis. Every
wedge of $e$ original forms lies generically in $M$ tensored by the
sum of their powers of $L$; saturation makes it a regular section of
that line. In the fixed basis, one coordinate of the unbounded family
is nonzero for unbounded $m$. Wedge those forms with the other $e-1$
fixed basis forms. These wedges are nonzero and invariant, and their
twist exponents are $m+1$ plus a fixed sum. They give the required
sections with $a_i\to\infty$.
\end{proof}


\begin{corollary}[Natural invariant nonvanishing]\label{thm:main}
Let $X$ be smooth connected projective over $\C$, with smoothly semipositive $L=-K_X$. If an algebraic torus $T$ acts on $X$ and $\chi(X,\mathcal O_X)\ne0$, then $H^0(X,mL)^T\ne0$ for some integer $m>0$, with the natural linearization. The torus may be trivial.
\end{corollary}
\begin{proof}
Every element of the connected torus is homotopic to the identity, so its action on singular cohomology is trivial. By Hodge decomposition its action on every $H^q(X,\mathcal O_X)$ is trivial. Consequently the invariant index at zero equals $\chi(X,\mathcal O_X)\ne0$.
Lemma~\ref{har:lem:forms} supplies either an invariant section directly or invariant sections of $M+a_iL$ with $a_i\to\infty$ and a nonzero map from $M$ into a positive tensor power of $\Omega_X^1$. The cotangent-tensor theorem \cite[Theorem~4.1]{LMPTX2023}, applied to the smooth zero-boundary pair with nef $-K_X$, says that $-M$ is pseudoeffective. Apply Theorem~\ref{inj:transfer} with $P=L$ naturally linearized.
\end{proof}

The harmonic constructions need additional information about the zero
divisors of these determinant sections. The following degree argument
provides it directly from the cotangent inclusion, independently of the
pseudoeffectivity theorem used in Corollary~\ref{thm:main}.

\begin{lemma}[Curvature-null determinant divisors]\label{har:lem:null-divisors}
Let $X$ be a smooth connected projective variety of dimension $n$, and
let $L=-K_X$ have smooth semipositive curvature $\theta$ of maximal
rank $r<n$. Suppose a line bundle $M$ admits a nonzero map
$M\to\bigwedge^e\Omega_X^p$ for integers $e,p>0$. Then
$c_1(M)G^{n-1}\le0$ for every Kähler class $G$. If
$0\ne s\in H^0(X,M+aL)$ for an integer $a$, its zero divisor $D$
satisfies
\[
 [D][\theta]^r[\omega]^{n-r-1}=0
\]
for every Kähler form $\omega$. Each component has zero intersection
separately, and $\theta^r$ vanishes on its smooth locus.
\end{lemma}

\begin{proof}
For every Kähler class \(G\), Yau's prescribed-Ricci theorem
\cite{Yau1978} gives a Kähler metric in \(G\) with Ricci form \(\theta\).
Kähler curvature symmetry identifies the mean curvature of \(T_X\)
with Ricci; the induced mean curvature of \(\bigwedge^e\Omega_X^p\)
is nonpositive. If \(\deg_G M>0\), solve a scalar Laplace
equation to give \(M\) a metric of positive constant mean curvature.
Its inverse tensor the cotangent construction then has negative mean
curvature, contradicting the Bochner formula for its nonzero section.
Thus \(c_1(M)G^{n-1}\le0\). Put \(G=[\theta]+t[\omega]\), divide by
\(t^{n-1-r}\), and let \(t\downarrow0\). Since \(\theta^{r+1}=0\),
\[
 c_1(M)[\theta]^r[\omega]^{n-1-r}\le0.
\]
The same intersection for \(M+aL\) is nonnegative by effectivity and
smooth semipositivity, and the contribution from $aL$ is zero. It is
consequently zero. Each component has zero intersection separately;
the nonnegative form $\theta^r\wedge\omega^{n-r-1}$ therefore vanishes
on its smooth locus, which gives the last assertion.
\end{proof}

\begin{remark}[The singular-norm proof of the degree sign]
The determinant degree calculation also has a useful current proof. In a K\"ahler metric with the prescribed semipositive Ricci form, let $e$ be the local image of a frame of $M$ in the cotangent construction. Its mean curvature is nonpositive. Chern's norm differentiation formula makes the trace of $i\partial\bar\partial\log(|e|^2+\varepsilon)$ nonnegative: the negative mean-curvature term has the required sign, and the derivative term is nonnegative by Cauchy--Schwarz. The logarithms converge locally in $L^1$, by comparison with logarithms of holomorphic coefficients. The resulting trace-positive current represents $-c_1(M)$. Integrating proves $c_1(M)G^{n-1}\le0$, including possible codimension-two zeros of the injected frame. This retains a singular-norm proof of the degree sign independently of the constant-mean-curvature contradiction used above.
\end{remark}
